\section{Limitations}\label{sec:limits}

\paragraph{Assumed technology constants.}
The engine electrical efficiency (0.43) and loss fraction (0.15 by default; 0.10 to 0.20 in the generators), the turbine efficiencies, the chiller reversible-limit fractions and the 0.8 COP cap, the heat-pump approach, the rejector approaches and fan fraction, the ORC fluid
vapor limits, the 333~K stack floor and the fraction of IT heat on the chilled loop are assumed (\cref{tab:constants}). They are not fitted to data, and the Design and Audit outcomes depend on them. In particular, the feasibility limits of \cref{sec:feasibility} follow from these constants.

\paragraph{Heat that is not recovered or modeled.}
Engine jacket and radiation heat is rejected, not recovered. The absorption chiller is a system-level model with a fraction of the reversible COP and a cap, not a LiBr property model.
Water use is an upper bound. At an 18\,$^\circ$C chilled-water supply the chiller sits at its COP cap, so the absorption results are insensitive to $\varepsilon$ there. Free cooling is not modeled, which makes the cold cells of \cref{fig:fig_design_feasibility} infeasible in the model.

\paragraph{Operation.}
There are no real part-load curves, no thermal storage, no dispatch, and no start-up or ramp limits. Trace results are bounds (\cref{sec:trace}), and the modular-fleet penalty depends on the trace shape.

\paragraph{No measured on-site-generation data center.}
Published measurements of a supercomputer set operating points, and nothing in this paper was measured on a data center with on-site generation. The Frontier trace evaluation is pending, and the only trace result is on a synthetic trace. There is no validation against a real plant: the cross-check in \cref{sec:env} is one operating point of one configuration against an independent implementation of the same equations.

\paragraph{Audit.}
The prime mover was the largest destroyer in every plant generated, so the ranking excludes it, and the task is a ranking among the remaining components. The best allowed change is mostly the pressure ratio or the turbine inlet temperature on turbine plants, where a fixed rule recovers \val{audb.rule_gain_id_pct}\% of the best gain when it is valid. The held-out plants (engine and steam-ORC) offer much smaller savings, so the gain channel is small in absolute terms there.

\paragraph{Ledger repair.}
The task is solvable by a rule: a program that rebuilds every line from the topology solves every row by construction. It measures whether a model applies that rule from text and localizes the fault, and it does not measure thermodynamic insight. Of the generated rows, \val{led.none_pct}\% show no second-law symptom, so the checks alone cannot detect them; the model has to find them from the text.

\paragraph{Repair.}
Most Repair faults are caught by parameter-range checks, and the baseline search solves almost all rows. The task is useful for the cost it measures (solver calls), not for difficulty of detection.

\paragraph{Reference and baselines.}
The Design reference is the best design found by enumeration plus a local polish, not a proven optimum, and the enumerate-then-polish curve runs the procedure that defines the reference. The baseline samples are small (\val{des.n_sites} Design sites, \val{aud.n_id} and \val{aud.n_ood} Audit plants, \val{rep.total_n} Repair rows, one example plant for each of the cross-check and the exergy-flow figure), and the intervals are wide where we give them.

\paragraph{Synthetic data.}
All training and evaluation data are sampled by the generators and labeled by the simulator and the verifiers. They are not measurements, and the distribution of plants is the one the generators produce.

\paragraph{Teacher and memory.}
A teacher answer that the tool suite accepts is correct for the checks that the suite has and may still be poor, for example a valid repair that is not the original plant; the checks decide validity and not optimality. Frontier models share habits with the student they teach, which narrows the variety we want, and different models differ in acceptance rate, so the corpus depends on which models were available within the quota and on their recorded prices. The memory is retrieved lexically, and notes mined from rejected or accepted attempts could leak information across splits unless the split of the mined records is controlled. Teacher-authored diagrams are checked by the same tool suite, so they inherit its blind spots: a plant that the verifiers cannot fault is accepted.

\paragraph{Diagram variety.}
The model's understanding of diagrams is tested on held-out diagrams of the catalog. The families in the catalog bound what this tests: a family that is absent cannot be held out, and held-out diagrams of a family that also appears in training test recombination more than novelty (\todo{limits of the final catalog}).

\paragraph{Model results and replay arms.}
At the time of writing no model has been trained or evaluated on a GPU, because none was available locally; the GPU code is reviewed but untested, and the known risks are loading Qwen3.5-4B through the causal-LM class, the TRL argument names, 4-bit loading and merging, and generation-length cutoffs. The runbook records them. Consequently every trained-model statement is pending. Frontier arms are evaluated by replaying saved answers: this gives an identical prompt and verifier for every arm, but it cannot give a frontier model solver feedback after the saved answer, so the frontier-with-tools arm is measured on the traces we collected during generation rather than by re-running the tools at evaluation time. We report this as a limitation of the comparison, not of the environment.

\paragraph{Verifier gaming and software status.}
Verifier gaming is possible. Edits are restricted to existing fields, the site check is in the Repair reward, and the reward functions are tested with junk and mutated answers for robustness (they must not raise and must stay bounded), but this does not exclude gaming by a trained policy. The project has \val{env.n_tests} tests that pass on the CPU; the continuous-integration workflow in the project has never been executed.

\paragraph{Summary of threats and mitigations.}
The main threats are construct validity (a verifier may accept a plant that a physicist would reject), internal validity (the baseline search may define the difficulty of Repair, and the reference design may define the difficulty of Design), external validity (all plants are generated from the same registry and no real plant is measured), and conclusion validity (small samples and many comparisons). We mitigate them, respectively, by publishing the verifier definitions and the accepted/rejected attempts, by reporting the procedures that define the baselines and the reference, by stating explicitly that transfer is between diagrams of the catalog and not to real plants, and by Wilson intervals, paired tests with Holm correction and fixed seeds. These mitigations bound the claims: the results support statements about this environment and its held-out diagrams, not about on-site generation in general.
